\begin{tabular}{llll}
\toprule
Family & Variable & Unit & Definition \\
\midrule
Baseline & Distance & m & Distance between transmission and reception (Haversine distance). \\
Baseline & Freq & MHz & Carrier frequency of the transmitted signal. \\
Baseline & Tx\_power & W & Effective isotropic radiated power at the transmitter. \\
Baseline & Tx\_height & m & Height of the transmitter antenna above ground level. \\
Baseline & Rx\_height & m & Height of the receiver antenna above ground level. \\
Terrain & Slope\_Tx\_Rx\_50m & dimensionless & Average slope between the transmission point and the point 50 m after reception along the Tx to Rx direction. \\
Terrain & Roughness\_Tx\_Rx\_50m & m & Rugosity (standard deviation of elevations) between the transmitter and the point 50 m downstream of the receiver along the Tx to Rx direction. \\
Geometry & Azimut\_Tx\_Rx & deg & Horizontal angle between true north and the Tx-Rx direction, measured clockwise in the range [0 deg, 360 deg). \\
Geometry & Tilt\_Tx\_Rx & deg & Elevation angle of the Tx-Rx line of sight. Positive: downtilt (receiver below transmitter); negative: uptilt (receiver above transmitter). \\
Obstacles & LOS & \{0,1\} & Binary indicator of line-of-sight between Tx and Rx. \\
Obstacles & First\_building\_m & m & Distance from Tx to the nearest building intercepting the Tx-Rx ray. \\
Obstacles & First\_tree\_m & m & Distance from Tx to the nearest tree intercepting the Tx-Rx ray. \\
Obstacles & Fresnel\_buildings & count & Number of buildings whose footprint intersects the first Fresnel ellipsoid of the Tx-Rx link. \\
Obstacles & Fresnel\_trees & count & Number of trees whose canopy intersects the first Fresnel ellipsoid of the Tx-Rx link. \\
Obstacles & Buildings\_near\_Rx & count & Number of buildings within a fixed radius (50 m) around the receiver. \\
Obstacles & Trees\_near\_Rx & count & Number of trees within a fixed radius (50 m) around the receiver. \\
Weather & Temp & C & Ambient air temperature at the measurement time. Retained (with Rhum) for its established link to atmospheric refractivity; Wdir/Wspd/Pres were dropped as they have no established physical role in point-to-area median-field prediction. \\
Weather & Rhum & \% & Relative humidity at the measurement time; drives the wet term of atmospheric refractivity together with Temp. \\
Context & Environment & cat. & Categorical morphology label: urban Cotonou vs. suburban Kandi. \\
\bottomrule
\end{tabular}
